%% Hand-authored circuitikz figure -- NOT generated by maestro2tex, placed by it via
%% the `order` list in configs/anatop_tb_ctrl_amp_ab_corners.json. Topology transcribed
%% from the Interactive.168 cell-drive bench netlist on 2026-08-16. Supplies, the bias
%% rails and En are omitted -- identical in every bench and carrying no measurement.
%% The cell is the abstracted symbol of sch_cell_symbol.tex -- its internal
%% network is not drawn anywhere in this chapter; the values are in the models
%% table.
\providecommand{\MaestroRoot}{include/analog/}
\input{\MaestroRoot sch_cell_symbol.tex}
\begin{figure}[htbp]
  \centering
  \resizebox{0.60\linewidth}{!}{%
\ctikzset{bipoles/length=1.0cm, resistors/scale=0.85, capacitors/scale=0.85, sources/scale=0.85}
\begin{circuitikz}[
    every node/.append style={font=\footnotesize},
    nlab/.style={font=\scriptsize, inner sep=1.5pt},
    opt/.style={font=\scriptsize, text=black!55, align=center, inner sep=2pt},
    railnode/.style={circle, fill=black, inner sep=1.05pt},
    prb/.style={draw, circle, minimum size=9pt, inner sep=0pt},
  ]
  \node[op amp] (A) at (0,0) {};
  \node at (A.center) {\footnotesize A1};
  %% three-electrode cell, drawn first so its stub ends exist for the wiring
  \CellSymThree{5.9,0}{0.85}
  \node[opt, anchor=north] at (5.9,-1.35) {three-electrode\\cell};
  %% Input pins leave horizontally; every branch that meets them is placed at the
  %% pin height with (x,0 |- pin) so no run is ever diagonal.
  \coordinate (Ap) at (A.+);
  \coordinate (Am) at (A.-);
  \coordinate (VREF) at (-1.7,0 |- Ap);
  \coordinate (FB)   at (-3.0,0 |- Am);
  \draw (A.+) -- (VREF);
  \draw (VREF) to[V, l_=$v_{\mathrm{ref}}$] (-1.7,-3.0) node[ground]{};
  \node[opt, anchor=north] at (-1.7,-3.45) {swept \SIrange{0}{2.5}{\volt}};
  %% reference-electrode feedback, through the loop-gain probe
  \node[prb] (P) at (-1.8,2.6) {};
  \draw (A.-) -- (FB) -- (-3.0,2.6) -- (P.west);
  \node[nlab, above=2pt] at (P.north) {probe};
  \draw (P.east) -- (5.9,2.6) -- (cs-re);
  %% external cascode compensation, C1 and C2 back onto the output
  \draw (A.up) -- (A.up |- 0,1.55);
  \draw (A.up |- 0,1.55) to[C, l=$C_c$] (2.2,1.55) -- (2.2,0);
  \node[nlab, anchor=east, xshift=-2pt] at (A.up |- 0,1.1) {\texttt{C1}};
  \draw (A.down) -- (A.down |- 0,-1.55);
  \draw (A.down |- 0,-1.55) to[C, l_=$C_c$] (2.2,-1.55) -- (2.2,0);
  \node[nlab, anchor=east, xshift=-2pt] at (A.down |- 0,-1.1) {\texttt{C2}};
  %% output node -> counter electrode
  \draw (A.out) -- (2.2,0);
  \node[railnode] at (2.2,0){};
  \draw (2.2,0) -- (cs-ce);
  \node[railnode] at (3.4,0){};
  \draw (3.4,0) to[C, l=$C_L$] (3.4,-3.0) node[ground]{};
  %% working electrode held at V_CM; the cell current is measured in that source
  \draw (cs-we) to[V, l=$V_{\mathrm{CM}}$] (7.23,-3.0) node[ground]{};
  \node[opt, anchor=north] at (7.23,-3.45) {cell current\\measured here};
\end{circuitikz}
  }
  \caption[{Cell-drive bench for the control amplifier.}]{Cell-drive bench. The amplifier drives the counter electrode of a three-electrode cell and closes its loop on the reference electrode, with the commanded potential $v_{\mathrm{ref}}$ swept over the full \SIrange{0}{2.5}{\volt} supply range. The working electrode is held at $V_{\mathrm{CM}}$ by an ideal source: no load current is injected, so the cell current is set by the sweep against the cell impedance and is measured in that source, and the regulation errors quoted at \SI{\pm10}{\micro\ampere} are read off the sweep at the points where the cell current reaches those values. The probe in the reference-electrode return measures the loop gain with the cell inside the loop; it is not the amplifier's own gain, and the two are tabulated separately. Cell elements are $R_{\mathrm{ce}} = R_u = \SI{1}{\kilo\ohm}$ and $R_{\mathrm{ct}} = \SI{10}{\kilo\ohm}$ in parallel with $C_{\mathrm{dl}} = \SI{1}{\micro\farad}$ (Figure~\ref{fig:randles-cell}), the load is $C_L = \SI{5}{\pico\farad}$, and each external compensation capacitor $C_c$ is \SI{500}{\femto\farad}.}
  \label{fig:tb-celldrive}
\end{figure}
